\chapter{Varietà differenziabili}
    \section{Varietà topologiche}
        \begin{definizione}
        \mbox{}

        Si definisce \textcolor{red}{varietà topologica} di dimensione $n \ge 1$, uno spazio topologico $M$ tale che:
        \begin{enumerate}[label=\arabic*)]
            \item $M$ è di Hausdorff (o $T_2$) 
            \textcolor{blu}{(cioè $\forall p, q \in M \quad \exists U, V \text{ intorni aperti di } p \text{ e } q \text{ risp. t.c. } U \cap V = \emptyset$)};
            \item $M$ è a base numerabile (o secondo-numerabile) \textcolor{blu}{(cioè $M$ ammette una base di topologia numerabile)};
            \item $M$ è localmente Euclideo di dimensione $n$ 
            \textcolor{blu}{(cioè $\forall p \in M \quad \exists U \text{ intorno aperto di } p, \ \exists \tilde{U} \subset \mathbb{R}^n \text{ aperto}, \ \exists \varphi : U \to \tilde{U} \text{ omeomorfismo}$)}
        \end{enumerate}
        \end{definizione}

        \begin{definizione}
        Sia $M$ varietà topologica di dimensione $n \ge 1$. Si definisce \textcolor{red}{carta} (o carta locale, o carta coordinata, o carta ammissibile) ogni coppia $(U, \varphi)$, con $U$ aperto di $M$ e $\varphi : U \to \varphi(U) = \tilde{U} \subset \mathbb{R}^n$ omeomorfismo tra $U$ e un aperto $\varphi(U)$ di $\mathbb{R}^n$.
        $U$ è detto \textcolor{red}{dominio della carta}, o \textcolor{red}{dominio coordinato}.

        Poniamo
        \[
        \forall q \in U \quad \varphi(q) = \left(x^1(q), \dots, x^n(q)\right) \in \varphi(U) \subset \mathbb{R}^n.
        \]
        Si definiscono quindi le funzioni
        \[
        x^i = p_i \circ \varphi : U \to \mathbb{R}
        \]
        dove $p_i : \varphi(U) \subset \mathbb{R}^n \to \mathbb{R}$, $(x^1, \dots, x^n) \mapsto x^i$ è la $i$-esima proiezione. Le funzioni $x^1, \dots, x^n : U \to \mathbb{R}$ sono dette funzioni \textcolor{red}{coordinate locali} su $U$ (o semplicemente coordinate locali) e sono funzioni continue.

        Dalla definizione di varietà topologica segue che $\forall p \in M \ \exists (U, \varphi)$ carta locale tale che $p \in U$. Pertanto, $M$ ammette una famiglia $\mathcal{A} = \{(U_\alpha, \varphi_\alpha)\}_{\alpha \in A}$ di carte locali tale che
        \[
        M = \bigcup_{\alpha \in A} U_\alpha \quad (\text{i domini delle carte ricoprono } M).
        \]
        La famiglia $\mathcal{A} = \{(U_\alpha, \varphi_\alpha)\}_{\alpha \in A}$ è detta \textcolor{red}{atlante}.
        \end{definizione}

        \begin{esempio}
        \begin{enumerate}[label=\arabic*)]
            \item $\mathbb{R}^n$ è varietà topologica di dimensione $n$. Un atlante è $\mathcal{A} = \{(\mathbb{R}^n, i_{\mathbb{R}^n})\}$.
            \item Se $\mathscr{C} \subset \mathbb{R}^n$ è una \textcolor{red}{curva differenziabile}, allora $\mathscr{C}$ è varietà topologica di dimensione $1$. Un atlante è $\{(U_\alpha, \varphi_\alpha)\}$ con
            \[
            \varphi_\alpha = \alpha^{-1} : U_\alpha \to I_\alpha \subset \mathbb{R}
            \]
            essendo $\alpha : I_\alpha \to U_\alpha \subset \mathscr{C} \subset \mathbb{R}^n$ omeomorfismo e parametrizzazione regolare.
            \item Se $S \subset \mathbb{R}^3$ è una \textcolor{red}{superficie differenziabile}, allora $S$ è varietà topologica di dimensione $2$. Un atlante è $\{(U_\alpha, \varphi_\alpha)\}$ con
            \[
            \varphi_\alpha = \alpha^{-1} : U_\alpha \to A_\alpha \subset \mathbb{R}^2
            \]
            essendo $\alpha : A_\alpha \to U_\alpha \subset S \subset \mathbb{R}^3$ omeomorfismo e parametrizzazione regolare.
            
            Pertanto sono varietà topologiche:
            \begin{itemize}
                \item $S^2 \subset \mathbb{R}^3$;
                \item il grafico di una funzione differenziabile $\varphi : A \subset \mathbb{R}^2 \to \mathbb{R}$ ($S = \{(u, v, \varphi(u,v)) \mid (u,v) \in A\}$);
                \item le superfici di livello di una funzione differenziabile $F : A \subset \mathbb{R}^3 \to \mathbb{R}$ ($S = F^{-1}(a)$, con $a$ valore regolare per $F$).
            \end{itemize}
        \end{enumerate}
        \end{esempio}
    \section{Funzione di transizione tra due carte}
        \begin{definizione}
        Sia $M$ varietà topologica di dimensione $n \ge 1$. Siano $(U_\alpha, \varphi_\alpha)$, $(U_\beta, \varphi_\beta)$ carte locali tali che $U_\alpha \cap U_\beta \neq \emptyset$. Si può considerare la seguente composizione:
        \[
        \varphi_\alpha(U_\alpha \cap U_\beta) \subset \mathbb{R}^n \xrightarrow{\varphi_\alpha^{-1}} U_\alpha \cap U_\beta \xrightarrow{\varphi_\beta} \varphi_\beta(U_\alpha \cap U_\beta) \subset \mathbb{R}^n.
        \]
        L'applicazione
        \[
        \varphi_\beta \circ \varphi_\alpha^{-1} : \varphi_\alpha(U_\alpha \cap U_\beta) \to \varphi_\beta(U_\alpha \cap U_\beta)
        \]
        è un omeomorfismo e prende il nome di \textcolor{red}{funzione di transizione} dalla carta $(U_\alpha, \varphi_\alpha)$ alla carta $(U_\beta, \varphi_\beta)$.

        Siano $x^1, \dots, x^n : U_\alpha \to \mathbb{R}$, $x^i = p_i \circ \varphi_\alpha$ le coordinate locali della carta $(U_\alpha, \varphi_\alpha)$ e $y^1, \dots, y^n : U_\beta \to \mathbb{R}$, $y^i = p_i \circ \varphi_\beta$ le coordinate locali della carta $(U_\beta, \varphi_\beta)$. Per ogni $p \in U_\alpha \cap U_\beta$:
        \[
        \varphi_\alpha(p) = \left(x^1(p), \dots, x^n(p)\right), \qquad \varphi_\beta(p) = \left(y^1(p), \dots, y^n(p)\right),
        \]
        \[
        \left(x^1(p), \dots, x^n(p)\right) \xmapsto{\varphi_\beta \circ \varphi_\alpha^{-1}} \left(y^1(p), \dots, y^n(p)\right).
        \]
        La funzione di transizione $\varphi_\beta \circ \varphi_\alpha^{-1}$, alle coordinate di un punto $p$ nella carta $(U_\alpha, \varphi_\alpha)$, associa le coordinate dello stesso punto nella carta $(U_\beta, \varphi_\beta)$.

        \begin{center}
        \begin{tikzpicture}[scale=0.95, >=Stealth]
            % Superficie M
            \draw[thick, cyan!60!black] (-3, 1.5) to[out=20, in=165] (2.5, 2.5) to[out=-65, in=55] (1.2, -0.5) to[out=140, in=0] (-3, 1.5);
            \node[cyan!60!black, font=\bfseries] at (-2.5, 2.1) {$M$};

            % Aperto U_alpha (giallo/oliva)
            \fill[olive!40, opacity=0.7] (-0.3, 1.8) to[out=10, in=150] (1.0, 2.0) to[out=-30, in=60] (0.8, 1.1) to[out=-120, in=-20] (-0.5, 1.3) to[out=160, in=-170] (-0.3, 1.8);
            \draw[olive!80!black, thick] (-0.3, 1.8) to[out=10, in=150] (1.0, 2.0) to[out=-30, in=60] (0.8, 1.1) to[out=-120, in=-20] (-0.5, 1.3) to[out=160, in=-170] (-0.3, 1.8);
            \node[olive!80!black] at (-0.4, 2.2) {$U_\alpha$};

            % Aperto U_beta (verde)
            \fill[green!50!black, opacity=0.45] (0.5, 1.95) to[out=20, in=130] (1.6, 1.7) to[out=-50, in=30] (1.4, 0.9) to[out=-150, in=-60] (0.3, 1.3) to[out=120, in=-160] (0.5, 1.95);
            \draw[green!50!black, thick] (0.5, 1.95) to[out=20, in=130] (1.6, 1.7) to[out=-50, in=30] (1.4, 0.9) to[out=-150, in=-60] (0.3, 1.3) to[out=120, in=-160] (0.5, 1.95);
            \node[green!40!black] at (1.9, 1.8) {$U_\beta$};

            % Intersezione U_alpha \cap U_beta (bordo rosso tratteggiato)
            \draw[red, thick, dashed] (0.5, 1.95) to[out=-160, in=120] (0.3, 1.3) to[out=-60, in=-120] (0.8, 1.1) to[out=60, in=-30] (1.0, 2.0) -- cycle;
            \node[red, font=\small] at (0.65, 0.8) {$U_\alpha \cap U_\beta$};

            % Punto p
            \fill[red] (0.65, 1.55) circle (1.5pt) node[below, font=\small] {$p$};

            % Aperto \varphi_\alpha(U_\alpha) a sinistra
            \draw[amber!90!black, very thick] (-3.2, -1.8) to[out=90, in=180] (-2.0, -0.9) to[out=0, in=90] (-1.0, -1.7) to[out=-90, in=0] (-2.2, -2.5) to[out=180, in=-90] (-3.2, -1.8);
            \node[amber!90!black, font=\small] at (-2.8, -0.6) {$\varphi_\alpha(U_\alpha) \subset \mathbb{R}^n$};

            % Immagine intersezione a sinistra
            \draw[red, thick] (-2.7, -1.6) to[out=60, in=170] (-1.8, -1.3) to[out=-10, in=80] (-1.5, -1.8) to[out=-100, in=-20] (-2.4, -2.0) to[out=160, in=-120] (-2.7, -1.6);
            \node[red, font=\scriptsize] at (-2.1, -1.15) {$\varphi_\alpha(U_\alpha \cap U_\beta)$};
            \fill[red] (-2.1, -1.7) circle (1.2pt) node[above, font=\scriptsize] {$\varphi_\alpha(p)$};

            % Aperto \varphi_\beta(U_\beta) a destra
            \draw[green!50!black, very thick] (2.4, -1.9) to[out=90, in=180] (3.6, -1.0) to[out=0, in=90] (4.6, -1.8) to[out=-90, in=0] (3.4, -2.7) to[out=180, in=-90] (2.4, -1.9);
            \node[green!50!black, font=\small] at (4.5, -0.7) {$\varphi_\beta(U_\beta) \subset \mathbb{R}^n$};

            % Immagine intersezione a destra
            \draw[red, thick] (2.9, -1.8) to[out=60, in=170] (3.8, -1.5) to[out=-10, in=80] (4.1, -2.0) to[out=-100, in=-20] (3.2, -2.2) to[out=160, in=-120] (2.9, -1.8);
            \node[red, font=\scriptsize] at (3.5, -1.35) {$\varphi_\beta(U_\alpha \cap U_\beta)$};
            \fill[red] (3.5, -1.9) circle (1.2pt) node[below, font=\scriptsize] {$\varphi_\beta(p)$};

            % Frecce delle carte
            \draw[->, thick] (-0.2, 1.3) to[out=-135, in=55] node[midway, above left] {$\varphi_\alpha$} (-1.8, -0.7);
            \draw[->, thin, gray!80!black] (-1.5, -0.8) to[out=50, in=-130] node[midway, below right, font=\small] {$\varphi_\alpha^{-1}$} (0.4, 1.2);
            \draw[->, thick] (1.6, 1.3) to[out=-35, in=110] node[midway, above right] {$\varphi_\beta$} (3.4, -0.8);

            % Freccia funzione di transizione
            \draw[->, red, thick] (-0.8, -1.8) to[out=-15, in=-165] node[midway, below] {$\varphi_\beta \circ \varphi_\alpha^{-1}$} (2.2, -1.8);
        \end{tikzpicture}
        \end{center}

        L'omeomorfismo inverso di $\varphi_\beta \circ \varphi_\alpha^{-1}$ è
        \[
        \varphi_\alpha \circ \varphi_\beta^{-1} : \varphi_\beta(U_\alpha \cap U_\beta) \to \varphi_\alpha(U_\alpha \cap U_\beta),
        \]
        funzione di transizione dalla carta $(U_\beta, \varphi_\beta)$ alla carta $(U_\alpha, \varphi_\alpha)$.
        \end{definizione}

        \begin{osservazione}
        Si consideri su una curva differenziabile $\mathscr{C} \subset \mathbb{R}^n$ un atlante differenziabile $\{(U_\alpha, \varphi_\alpha)\}$ dove $\varphi_\alpha = \alpha^{-1}$, con $\alpha : I_\alpha \subset \mathbb{R} \to U_\alpha \subset \mathscr{C} \subset \mathbb{R}^n$ parametrizzazione regolare e omeomorfismo. La funzione di transizione
        \[
        \varphi_\beta \circ \varphi_\alpha^{-1} = \beta^{-1} \circ \alpha : I_\alpha \to I_\beta
        \]
        è il cambiamento di parametro che (è stato provato) è un diffeomorfismo.

        Analogamente per una superficie differenziabile $S \subset \mathbb{R}^3$, considerato un atlante differenziabile $\{(U_\alpha, \varphi_\alpha)\}$ tale che $\varphi_\alpha = \alpha^{-1}$ con $\alpha : A_\alpha \subset \mathbb{R}^2 \to U_\alpha \subset S \subset \mathbb{R}^3$ parametrizzazione regolare e omeomorfismo, la funzione di transizione
        \[
        \varphi_\beta \circ \varphi_\alpha^{-1} = \beta^{-1} \circ \alpha : A_\alpha \to A_\beta
        \]
        è il cambiamento di parametri ed è un diffeomorfismo.
        \end{osservazione}
        \newpage
        \begin{proposizione}
        Sia $M$ varietà topologica e sia $p \in M$. I domini delle carte locali costituiscono un \textcolor{red}{sistema fondamentale di intorni} di $p$, ovvero
        \[
        \forall A \text{ intorno aperto di } p \quad \exists (U, \varphi) \text{ carta locale t.c. } p \in U \subset A.
        \]
        \end{proposizione}

        \begin{proof}
        Sia $A$ intorno aperto di $p$. Considerato $p \in M$, $\exists (U', \varphi')$ carta locale tale che $p \in U'$.
        Poniamo $U = U' \cap A$ intorno aperto di $p$.
        Allora $\varphi = (\varphi'|_{U})_* : U \to \varphi(U) \subset \mathbb{R}^n$ è omeomorfismo e $(U, \varphi)$ è carta locale in $p$ con $U \subset A$.
        \end{proof}

        \begin{proposizione}
        Ogni varietà topologica è \textcolor{red}{localmente connessa}, ovvero ogni punto ammette un sistema fondamentale di intorni connessi, ovvero
        \[
        \forall p \in M \ \forall A \text{ intorno aperto di } p \quad \exists V \text{ intorno connesso di } p \text{ tale che } V \subset A.
        \]
        \end{proposizione}

        \begin{proof}
        Sia $p \in M$, sia $A$ intorno aperto di $p$. Sia $(U, \varphi)$ carta locale tale che $p \in U \subset A$.
        Poiché $\varphi : U \to \varphi(U) \subset \mathbb{R}^n$ è omeomorfismo, sia $B = B_r(\varphi(p)) \subset \varphi(U)$.
        $B$ è connesso di $\mathbb{R}^n$ e quindi di $\varphi(U)$.
        Allora $p \in \varphi^{-1}(B) \subset U \subset A$ e $V := \varphi^{-1}(B)$ è intorno connesso di $p$.
        \end{proof}

        \begin{proposizione}
        Sia $M$ varietà topologica. Ogni punto di $M$ ammette un \textcolor{red}{sistema fondamentale di intorni compatti}, ovvero
        \[
        \forall p \in M \ \forall A \text{ intorno aperto di } p \quad \exists K \text{ intorno compatto di } p \text{ tale che } K \subset A.
        \]
        \end{proposizione}

        \begin{proof}
        Sia $p \in M$ e sia $A$ intorno aperto di $p$. Sia $(U, \varphi)$ carta locale in $p$ t.c. $U \subset A$.
        Poiché $\varphi : U \to \varphi(U) \subset \mathbb{R}^n$ è omeomorfismo, sia $\overline{B} = \overline{B_r(\varphi(p))} \subset \varphi(U)$.
        $\overline{B}$ è intorno compatto di $\varphi(p)$.
        Allora $K := \varphi^{-1}(\overline{B})$ è intorno compatto di $p$ tale che $K \subset U \subset A$.
        \end{proof}

        \begin{osservazione}
        In particolare, una varietà topologica è \textcolor{red}{localmente compatta} (ogni punto possiede un intorno compatto).
        \end{osservazione}
        \newpage
        \begin{definizione}
        Sia $M$ varietà topologica, $\dim M = n \ge 1$. Siano $(U, \varphi)$, $(U', \varphi')$ carte locali. Si dice che $(U, \varphi)$ e $(U', \varphi')$ sono \textcolor{red}{compatibili} (o differenziabilmente compatibili) se:
        \begin{itemize}
            \item $U \cap U' = \emptyset$, oppure
            \item $U \cap U' \neq \emptyset$ e la funzione di transizione
            \[
            \varphi' \circ \varphi^{-1} : \varphi(U \cap U') \to \varphi'(U \cap U')
            \]
            è un \textcolor{red}{diffeomorfismo} (di classe $C^\infty$).
        \end{itemize}
        \end{definizione}

        \begin{osservazione}
        Si possono studiare anche varietà differenziabili per le quali si richiede che le funzioni di transizione siano differenziabili di classe $C^k$, o anche che siano analitiche reali (ovvero esprimibili come una serie convergente di potenze).
        \end{osservazione}

        \begin{definizione}
        Sia $M$ varietà topologica, $\dim M = n$. Si definisce \textcolor{red}{atlante differenziabile} su $M$ un atlante $\mathcal{A} = \{(U_\alpha, \varphi_\alpha)\}_{\alpha \in A}$ le cui carte locali siano a due a due differenziabilmente compatibili.
        Pertanto, un atlante differenziabile è una collezione di carte $\mathcal{A} = \{(U_\alpha, \varphi_\alpha)\}_{\alpha \in A}$ tale che:
        \begin{itemize}
            \item $M = \bigcup_{\alpha \in A} U_\alpha$;
            \item $\forall \alpha \in A \quad \varphi_\alpha : U_\alpha \to \varphi_\alpha(U_\alpha) \subset \mathbb{R}^n$ è \textcolor{red}{omeomorfismo} tra $U_\alpha$ aperto di $M$ e $\varphi_\alpha(U_\alpha)$ aperto di $\mathbb{R}^n$;
            \item $\forall \alpha, \beta \in A$ t.c. $U_\alpha \cap U_\beta \neq \emptyset$, l'applicazione
            \[
            \varphi_\beta \circ \varphi_\alpha^{-1} : \varphi_\alpha(U_\alpha \cap U_\beta) \to \varphi_\beta(U_\alpha \cap U_\beta)
            \]
            è un \textcolor{red}{diffeomorfismo}.
        \end{itemize}
        \end{definizione}

    \section{Varietà differenziabili}
    \section{Esempi di varietà differenziabili}
    \section{Funzioni differenziabili su una varietà differenziabile}
    \section{Applicazioni tra varietà differenziabili}
    \section{Proprietà topologiche delle varietà differenziabili}
    \section{Partizione dell'unità}
    \section{Funzioni bump su varietà differenziabili}